% =============================================================
%  Chapter 5 — Synthesis: A Hierarchy of Complexity in
%              Continual Hamiltonian Learning
% =============================================================

\chapter{Synthesis: A Hierarchy of Complexity in Continual Hamiltonian Learning}
\label{chap:hierarchy}

\providecommand{\Ham}{\mathcal{H}}
\providecommand{\Potential}{\mathcal{R}}
\providecommand{\PhaseState}{z}
\providecommand{\Pos}{q}
\providecommand{\Mom}{p}
\providecommand{\Mass}{\mathrm{M}}
\providecommand{\J}{\mathbb{J}}

% -----------------------------------------------------------------------
\section{Why This Chapter Is Needed}
\label{sec:ch5:why}
% -----------------------------------------------------------------------

Chapters~\ref{chap:grlsnam} and~\ref{chap:material} should not be read
as two independent navigation systems.
They are two instances of the same structured learner, operating at
different levels of richness.
The central contribution of the thesis is not only a better navigator
in either setting: it is the \textbf{structured progression} itself—the
fact that the learner can be enlarged in a principled, coupled way
without discarding earlier structure.
x₹₹x₹₹
This chapter makes that progression explicit.
It answers four questions:
\begin{enumerate}
  \item What shared skeleton underlies both chapters?
  \item What changed from Chapter~\ref{chap:grlsnam} to
        Chapter~\ref{chap:material}, and why were those changes not
        superficial?
  \item What does the enrichment principle reveal about a general
        framework for continual Hamiltonian navigation?
  \item What remains unsolved, and what does the next stage look like?
\end{enumerate}

The answers motivate the broader thesis claim:

\begin{quote}\itshape
The thesis develops a structured Hamiltonian learning framework in
which geometry, field topology, adaptation timescale, and stochastic objective
richness are not separate design choices, but \textbf{coupled layers of
a single enrichable phase-space learner}.
\end{quote}

% -----------------------------------------------------------------------
\section{The Common Skeleton}
\label{sec:ch5:skeleton}
% -----------------------------------------------------------------------

Despite their different energy landscapes, objectives, and adaptation
loops, Chapters~\ref{chap:grlsnam} and~\ref{chap:material} share an
identical structural template.

\begin{definition}[Continual Hamiltonian Navigation Skeleton]
\label{def:ch5:skeleton}
A \emph{continual Hamiltonian navigation instance} is a tuple
\begin{equation}
  \mathfrak{S} = (H_\theta,\, R_\theta,\, f_\xi,\, \mathcal{L},\,
                  \mathcal{U}),
  \label{eq:ch5:skeleton_tuple}
\end{equation}
where:
\begin{itemize}
  \item $H_\theta : T^*\mathcal{Q} \to \mathbb{R}$ is a parameterized
        stored energy with coefficient vector $\theta$;
  \item $R_\theta : T^*\mathcal{Q} \to T^*_\Pos\mathcal{Q}$ is the
        dissipative port (separate from $H_\theta$);
  \item $f_\xi : \mathcal{E} \to \theta$ is the meta-navigator mapping
        local context to energy coefficients;
  \item $\mathcal{L}$ is a multi-term objective over rollout trajectories;
  \item $\mathcal{U}$ is the multi-timescale update hierarchy.
\end{itemize}
The \emph{policy} is the flow map
$\pi = \Phi^{\Delta t}_{H_\theta - R_\theta}$,
and the phase-space dynamics are
\begin{equation}
  \PhaseState_{t+1}
  = \underbrace{\Phi^{\tau}_{H_\theta}(\PhaseState_t)}_{\text{conservative}}
  - \tau\,R_\theta(\PhaseState_t)
  \;=:\; \Phi_{\vartheta_t}(\PhaseState_t,\, y_t),
  \label{eq:ch5:skeleton_dynamics}
\end{equation}
where $\vartheta_t = f_\xi(y_t, \text{goal})$ are the instantaneous
energy coefficients and $y_t$ is the local context observation.
\end{definition}

The shared template is:
\begin{equation}
  \boxed{
  \PhaseState_{t+1} = \Phi_{\vartheta_t}(\PhaseState_t, y_t),
  \qquad
  \vartheta_t = f_\xi(y_t, \vec{g}).
  }
  \label{eq:ch5:template}
\end{equation}
What changes between instances is \emph{what lives inside $H_\theta$,
$R_\theta$, $\mathcal{L}$, and $\mathcal{U}$}—the components of
$\mathfrak{S}$—not the skeleton itself.

Table~\ref{tab:ch5:instances} shows how the two thesis chapters
instantiate this skeleton.

\begin{table}[t]
\centering
\caption{%
  \textbf{Two instantiations of the continual Hamiltonian navigation
  skeleton.}
  The shared template \eqref{eq:ch5:template} is identical in both
  chapters; all variation is in the contents of $H_\theta$,
  $R_\theta$, $\mathcal{L}$, and $\mathcal{U}$.
}
\label{tab:ch5:instances}
\small
\setlength{\tabcolsep}{5pt}
\renewcommand{\arraystretch}{1.35}
\resizebox{\textwidth}{!}{%
\begin{tabular}{lll}
\toprule
\textbf{Component} &
\textbf{Stage 1: GRL-SNAM (Ch.~\ref{chap:grlsnam})} &
\textbf{Stage 2: Material-Aware (Ch.~\ref{chap:material})} \\
\midrule
$H_\theta$ &
$H_{\mathrm{kin}} + H_{\mathrm{geom}}$ &
$H_{\mathrm{kin}} + H_{\mathrm{geom}} + H_{\mathrm{mat}}^{\mathrm{soft}}
 + H_{\mathrm{mat}}^{\mathrm{hard}}$ \\
$R_\theta$ &
$\gamma\Mass^{-1}\Mom$ (implicit) &
$\gamma\Mass^{-1}\Mom$ (explicit separate port) \\
Context $y_t$ &
$\{d_j, \vec{n}_j\}$ (geometry only) &
$\{d_j, \vec{n}_j\} \cup \{\tilde{r}, \phi, P_t\}$ \\
$\theta$ &
$(\beta, \{\alpha_j\}, \gamma)$ &
$(\beta, \{\alpha_j\}, \gamma, \lambda_s, \lambda_h)$ \\
$\mathcal{L}$ &
$\mathcal{L}_{\mathrm{geom}}$ (expected cost) &
$\mathcal{L}_{\mathrm{nav}} + \mathcal{R}_H + \mathcal{R}_{\nabla H}$
  (CVaR + double regularization) \\
$\mathcal{U}$ &
Online $\alpha_j$ update (1 timescale) &
Segment $\to$ Episode $\to$ Curriculum (3 timescales) \\
Force channels &
$F_{\mathrm{goal}},\, F_{\mathrm{bar}},\, -R_\theta$ &
$F_{\mathrm{goal}},\, F_{\mathrm{bar}},\, F_{\mathrm{soft}},\,
 F_{\mathrm{hard}},\, -R_\theta$ \\
\bottomrule
\end{tabular}}
\end{table}

% -----------------------------------------------------------------------
\section{The Enrichment Principle}
\label{sec:ch5:principle}
% -----------------------------------------------------------------------

\subsection{Formal Statement}
\label{subsec:ch5:principle_formal}

The central structural observation of the thesis is:

\begin{theorem}[Enrichment Principle]
\label{thm:ch5:enrichment}
Let $\mathfrak{S}^{(k)}$ be a continual Hamiltonian navigation
instance at complexity level $k$.
Suppose a new energy term $H_{\mathrm{new}}(\Pos;\,\mathcal{E})$ is
added to $H_\theta^{(k)}$, where $\mathcal{E}$ contains new context
information not present at level $k$.
Then the following chain of consequences is forced:

\begin{enumerate}[(i)]
  \item \textbf{New force channel.}
        $F_{\mathrm{new}} = -\nabla_\Pos H_{\mathrm{new}}$ contributes
        to $F_{\mathrm{tot}}$, modifying the dynamics at every
        integration step.

  \item \textbf{New coefficient subspace.}
        $H_{\mathrm{new}}$ introduces at least one new learnable
        coefficient $\vartheta_{\mathrm{new}}$ into $f_\xi$'s output
        space.

  \item \textbf{New objective term.}
        The task criterion must be extended to measure performance
        along the new force channel, adding at least one term to
        $\mathcal{L}$.

  \item \textbf{New gradient pathway.}
        The rollout sensitivity recursion
        \eqref{eq:ch4:dp_dtheta}--\eqref{eq:ch4:dq_dtheta}
        acquires a new additive term
        $\partial\nabla_\Pos H_{\mathrm{new}}/\partial\theta$
        at each step.

  \item \textbf{New update law.}
        The new gradient pathway may require a new or modified
        adaptation rule in $\mathcal{U}^{(k+1)}$, since the
        coefficient $\vartheta_{\mathrm{new}}$ may operate at a
        different timescale from existing coefficients.
\end{enumerate}

The resulting enriched instance $\mathfrak{S}^{(k+1)}$ has the same
skeleton \eqref{eq:ch5:template} as $\mathfrak{S}^{(k)}$, with
strictly larger $H_\theta$, $\theta$, $\mathcal{L}$, and $\mathcal{U}$.
\end{theorem}

\begin{proof}[Proof sketch]
(i) follows directly from the definition of the potential energy's
gradient in Hamilton's equations.
(ii) follows because $H_{\mathrm{new}}$ introduces at least one weight
that modulates its contribution; if this weight were fixed it could be
absorbed into $H_{\mathrm{new}}$ directly, yielding no new degree of
freedom.
(iii) follows because without a loss term sensitive to
$F_{\mathrm{new}}$, the optimizer has no gradient signal to set
$\vartheta_{\mathrm{new}}$ away from zero.
(iv) follows by direct differentiation of \eqref{eq:ch4:p_update}.
(v) follows because $\vartheta_{\mathrm{new}}$ may interact with
the dynamics at a finer or coarser timescale than existing
coefficients, requiring a new loop in $\mathcal{U}$.
\end{proof}

\paragraph{Compact form.}
The enrichment principle can be written as a chain:
\begin{equation}
  H_{\mathrm{new}}
  \;\xrightarrow{(i)}\; F_{\mathrm{new}}
  \;\xrightarrow{(ii)}\; \vartheta_{\mathrm{new}}
  \;\xrightarrow{(iii)}\; \mathcal{L}_{\mathrm{new}}
  \;\xrightarrow{(iv)}\; \nabla_{\vartheta_{\mathrm{new}}}\mathcal{L}
  \;\xrightarrow{(v)}\; \text{update law}_{\mathrm{new}}.
  \label{eq:ch5:chain}
\end{equation}
Every step in this chain is a \emph{necessity}, not a design choice.
This is what distinguishes enrichment from simple extension: adding
$H_{\mathrm{new}}$ to the Hamiltonian obligates all five downstream
consequences.

\subsection{Tracing the Chain through the Thesis}
\label{subsec:ch5:chain_traced}

Table~\ref{tab:ch5:chain} traces \eqref{eq:ch5:chain} for each energy
term added in the thesis, making the coupling between all five links
explicit.

\begin{table}[t]
\centering
\caption{%
  \textbf{Enrichment chain for each energy term across both thesis stages.}
  Each row reads from left to right as the forced consequence of
  adding one energy term.
}
\label{tab:ch5:chain}
\small
\setlength{\tabcolsep}{4pt}
\renewcommand{\arraystretch}{1.4}
\begin{tabular}{lllll}
\toprule
\makecell{\textbf{Energy term}\\\textbf{(i) force}} &
\makecell{\textbf{(ii) coeff.}} &
\makecell{\textbf{(iii) objective term}} &
\makecell{\textbf{(iv) gradient}\\\textbf{pathway}} &
\makecell{\textbf{(v) update law}} \\
\midrule
\makecell[l]{$H_{\mathrm{geom}}$:\\ $F_{\mathrm{goal}},F_{\mathrm{bar}}$}
  & $\beta, \{\alpha_j\}$
  & $\mathcal{L}_{\mathrm{traj}}, \mathcal{L}_{\mathrm{clear}}$
  & $\partial\nabla D_\mathcal{G}/\partial\beta$
  & Online $\alpha_j$ (1 scale) \\
\midrule
\makecell[l]{$H_{\mathrm{mat}}^{\mathrm{soft}}$:\\ $F_{\mathrm{soft}}=-\lambda_s\nabla\tilde{r}$}
  & $\lambda_s$
  & $w_r\sum_t\tilde{r}(\vec{c}_t)\Delta s_t$
  & $\partial(-\lambda_s\nabla\tilde{r})/\partial\lambda_s$
  & CVaR episode grad. \\
\makecell[l]{$H_{\mathrm{mat}}^{\mathrm{hard}}$:\\ $F_{\mathrm{hard}}=-\lambda_h\nabla b(\phi)$}
  & $\lambda_h$
  & $w_h\sum_t\mathbf{1}[\phi<1\,\text{m}]$
  & $\partial(-\lambda_h\nabla b)/\partial\lambda_h$
  & CVaR episode grad. \\
\makecell[l]{$R_\theta$ port:\\ $-\gamma\Mass^{-1}\Mom$}
  & $\gamma$
  & $w_\gamma|\gamma-\gamma_0|^2$
  & $\partial(-\gamma\Mass^{-1}\Mom)/\partial\gamma$
  & Friction alignment \\
\makecell[l]{$\mathcal{R}_{\nabla H}$:\\ (field regularizer)}
  & —
  & $w_c\mathcal{L}_{\mathrm{clear}}+w_m\mathcal{L}_{\mathrm{multi}}$
  & $\partial F_{\mathrm{tot}}/\partial\theta$
  & Curriculum pass. \\
\bottomrule
\end{tabular}
\end{table}

Figure~\ref{fig:ch5:enrichment_ladder} gives a visual summary of the
full enrichment progression, from the geometry-only base to the richer
material-aware stage and beyond.

\begin{figure}[htbp]
\centering
\begin{tikzpicture}[
  font=\small\sffamily,
  % Stage box style
  stage/.style={
    rectangle, rounded corners=6pt,
    draw=#1!70!black, fill=#1!9,
    line width=1.4pt,
    minimum width=13.5cm, minimum height=2.0cm,
    align=left, inner sep=10pt
  },
  stageno/.style={
    circle,
    draw=#1!80!black, fill=#1!25,
    line width=1.2pt,
    minimum size=26pt,
    font=\bfseries\normalsize\sffamily
  },
  % Chain node style
  chain/.style={
    rectangle, rounded corners=3pt,
    draw=#1!60!black, fill=#1!12,
    line width=0.8pt,
    minimum width=1.9cm, minimum height=0.7cm,
    align=center, inner sep=3pt,
    font=\scriptsize\sffamily
  },
  % Arrow styles
  harrow/.style={-{Stealth[length=5pt]}, line width=1pt, color=gray!60},
  varrow/.style={-{Stealth[length=7pt]}, line width=1.4pt, color=#1},
  addlbl/.style={font=\scriptsize\sffamily\itshape, text=gray!60!black},
  % Dashed future
  futurebox/.style={
    rectangle, rounded corners=6pt,
    draw=gray!40, fill=gray!5,
    line width=1pt, dashed,
    minimum width=13.5cm, minimum height=1.6cm,
    align=left, inner sep=10pt
  },
]

% ── Stage 0: Geometry-only (Ch 3) ─────────────────────────────────────
\node[stage=teal] (s0) at (0,0) {%
  \hspace{30pt}%
  \begin{minipage}{12.2cm}
    \textbf{Stage 1 — Geometry-Only (Chapter~\ref{chap:grlsnam})}\\[3pt]
    $H^{(1)} = H_{\mathrm{kin}} + \beta D_\mathcal{G}
               + \sum_j\alpha_j b(d_j)$,\quad
    $R^{(1)}: \gamma\Mass^{-1}\Mom$ (implicit)\\[2pt]
    $\mathcal{L}^{(1)} = \mathcal{L}_{\mathrm{geom}}$,\quad
    $\mathcal{U}^{(1)}$: 1-scale online $\alpha_j$ update\\[2pt]
    \textit{Solves:} geometry-aware navigation under local sensing.
    \textit{Cannot:} reason about latent material risk.
  \end{minipage}
};
\node[stageno=teal] at ($(s0.west)+(18pt,0)$) {\textcolor{teal!80!black}{1}};

% ── Transition annotation ────────────────────────────────────────────
\node[addlbl] at ($(s0.south)!0.5!(s1.north)+(0,-0.18)$) {};

% % ── Enrichment chain (floating between stages) ─────────────────────
% \node[chain=orange] (cH)  at (-5.2,-2.05) {$+H_{\mathrm{mat}}$};
% \node[chain=orange] (cF)  at (-2.7,-2.05) {$\to F_{\mathrm{mat}}$};
% \node[chain=orange] (cV)  at (-0.2,-2.05) {$\to \vartheta_{\mathrm{new}}$};
% \node[chain=orange] (cL)  at ( 2.3,-2.05) {$\to \mathcal{L}_{\mathrm{new}}$};
% \node[chain=orange] (cU)  at ( 4.8,-2.05) {$\to \text{update}$};
% \draw[harrow] (cH.east) -- (cF.west);
% \draw[harrow] (cF.east) -- (cV.west);
% \draw[harrow] (cV.east) -- (cL.west);
% \draw[harrow] (cL.east) -- (cU.west);
% \node[addlbl, above=2pt of cH] {\textbf{Enrichment chain} (Theorem~\ref{thm:ch5:enrichment})};

% ── Stage 1: Material-aware (Ch 4) ──────────────────────────────────
\node[stage=orange, below=1.5cm of s0] (s1) {%
  \hspace{30pt}%
  \begin{minipage}{12.2cm}
    \textbf{Stage 2 — Material-Aware (Chapter~\ref{chap:material})}\\[3pt]
    $H^{(2)} = H^{(1)} + \lambda_s\tilde{r}(\vec{c})
               + \lambda_h b(\phi(\vec{c}))$,\quad
    $R^{(2)}: \gamma\Mass^{-1}\Mom$ (explicit port)\\[2pt]
    $\mathcal{L}^{(2)} = \mathrm{CVaR}_{0.95}(J) +
               \mathcal{R}_H + \mathcal{R}_{\nabla H}$,\quad
    $\mathcal{U}^{(2)}$: segment $\to$ episode $\to$ curriculum\\[2pt]
    \textit{Solves:} risk-sensitive navigation under oracle material.
    \textit{Cannot:} infer material from raw perception.
  \end{minipage}
};
\node[stageno=orange] at ($(s1.west)+(16pt,0)$)
  {\textcolor{orange!80!black}{2}};

\draw[varrow=orange!70!black] (s0.south) -- (s1.north);

% % ── Enrichment chain (between s1 and s2_future) ─────────────────────
% \node[chain=gray] (dH) at (-5.2,-4.1) {\small$+H_{\mathrm{perc}}$};
% \node[chain=gray] (dF) at (-2.7,-4.1) {\small$\to F_{\mathrm{perc}}$};
% \node[chain=gray] (dV) at (-0.2,-4.1) {\small$\to \vartheta_{\mathrm{perc}}$};
% \node[chain=gray] (dL) at ( 2.3,-4.1) {\small$\to \mathcal{L}_{\mathrm{perc}}$};
% \node[chain=gray] (dU) at ( 4.8,-4.1) {\small$\to \text{update}$};
% \draw[harrow] (dH.east) -- (dF.west);
% \draw[harrow] (dF.east) -- (dV.west);
% \draw[harrow] (dV.east) -- (dL.west);
% \draw[harrow] (dL.east) -- (dU.west);
% \node[addlbl, above=2pt of dH, text=gray!60]
%   {\textit{next enrichment (future)}};

% ── Stage 2: Future / deferred ──────────────────────────────────────
\node[futurebox, below=1.5cm of s1] (s2) {%
  \hspace{30pt}%
  \begin{minipage}{12.2cm}
    \textbf{Stage 3 — Perception-Led Material Inference (Future Work)}\\[3pt]
    $H^{(3)} = H^{(2)} + H_{\mathrm{perc}}(\vec{c}; \hat\rho)$,\quad
    learned material belief $\hat\rho$ from RGB-IR\\[2pt]
    $\mathcal{L}^{(3)}$ adds perception loss, uncertainty calibration;
    $\mathcal{U}^{(3)}$ adds belief update loop\\[2pt]
    \textit{Would solve:} end-to-end perception-to-navigation under
    latent material uncertainty.
  \end{minipage}
};
\node[stageno=gray!60] at ($(s2.west)+(16pt,0)$)
  {\textcolor{gray!50}{\textbf{3}}};
\draw[varrow=gray!50, dashed] (s1.south) -- (s2.north);

% ── Right-side axis label ──────────────────────────────────────────
\draw[{Stealth[length=6pt]}-{Stealth[length=6pt]},
      line width=1pt, color=gray!50]
  ($(s0.north east)+(0.5cm,0)$) -- ($(s2.south east)+(0.5cm,0)$)
  node[midway, right=6pt, font=\small\sffamily\bfseries,
       text=gray!60, rotate=-90, anchor=south] {enrichment axis};

\end{tikzpicture}
\caption{%
  \textbf{Enrichment ladder.}
  Each stage adds new energy terms, force channels, coefficients,
  objectives, and update laws.
  The transition arrow is labeled with the enrichment chain
  (Theorem~\ref{thm:ch5:enrichment}).
  Stage~3 (dashed) is future work: perception-led material inference
  via a learned belief term $H_{\mathrm{perc}}$.
}
\label{fig:ch5:enrichment_ladder}
\end{figure}

% -----------------------------------------------------------------------
\section{The Four Coupled Enlargements in Detail}
\label{sec:ch5:four}
% -----------------------------------------------------------------------

\subsection{Enlargement A: Stored Energy}
\label{subsec:ch5:energyA}

At Stage~1 the stored energy is:
\begin{equation}
  H^{(1)}(\Pos,\Mom)
  = \tfrac{1}{2}\Mom^\top\Mass^{-1}\Mom
    + \beta D_\mathcal{G}(\Pos,\vec{g})
    + \textstyle\sum_j \alpha_j b(d_j(\Pos)).
  \label{eq:ch5:H1}
\end{equation}
At Stage~2 it is:
\begin{equation}
  H^{(2)}(\Pos,\Mom)
  = H^{(1)} + \lambda_s\,\tilde{r}(\vec{c})
              + \lambda_h\,b_{\mathrm{sp}}(\phi(\vec{c})).
  \label{eq:ch5:H2}
\end{equation}
The geometry sub-Hamiltonian $H_{\mathrm{geom}}$ is carried forward
\emph{unchanged}; the material terms are \emph{additive}, not
replacements.
This is the direct mathematical embodiment of "enlargement rather than
replacement."

\begin{remark}[Sign of the potential]
\label{rem:ch5:sign}
In both stages the potential $\Potential = H - H_{\mathrm{kin}}$
is defined with the sign convention from the PMP reduction
(§\ref{subsec:prelim:pmp}): the gradient
$-\nabla_\Pos\Potential$ provides the restoring force driving
goal attraction and barrier repulsion.
The material terms enter with the same sign convention:
$-\lambda_s\nabla\tilde{r}$ pushes the robot away from high-risk
terrain, while $-\lambda_h\nabla b(\phi)$ repels it from hard hazard
boundaries.
\end{remark}

\subsection{Enlargement B: Coefficient and Control Space}
\label{subsec:ch5:coeffB}

The coefficient space grows from $\theta^{(1)} = (\beta, \{\alpha_j\}, \gamma)$
to $\theta^{(2)} = (\beta, \{\alpha_j\}, \gamma, \lambda_s, \lambda_h)$.
Two new dimensions are added; they are orthogonal to the existing
geometry coefficients in the sense that setting $\lambda_s = \lambda_h = 0$
exactly recovers Stage~1 behavior.

\begin{definition}[Coefficient subspace compatibility]
\label{def:ch5:compat}
An enrichment from $\theta^{(k)}$ to $\theta^{(k+1)}$ is
\emph{backward-compatible} if there exist values of the new
coefficients $\vartheta_{\mathrm{new}}$ such that
$H^{(k+1)}(\cdot;\theta^{(k+1)}\big|_{\vartheta_{\mathrm{new}}=0})
= H^{(k)}(\cdot;\theta^{(k)})$.
Both thesis transitions satisfy this property: the new coefficients
default to zero, recovering the lower-stage dynamics.
\end{definition}

Backward compatibility is important for the Phase~1 $\to$ Phase~2
curriculum (Algorithm~\ref{algo:ch4:material}): initializing from
Phase~1 weights with $\lambda_s = \lambda_h = 0$ provides a warm
start, and the curriculum only unfreezes these coefficients after the
geometry backbone has converged.

\subsection{Enlargement C: Objective and Risk Sensitivity}
\label{subsec:ch5:objectiveC}

The objective enlarges from expected-cost geometry supervision to a
CVaR-based multi-term criterion:
\begin{align}
  \mathcal{L}^{(1)}
  &= w_t\mathcal{L}_{\mathrm{traj}} + w_v\mathcal{L}_{\mathrm{vel}}
     + w_\gamma|\gamma-\gamma_0|^2 + w_c\mathcal{L}_{\mathrm{clear}}
     + w_m\mathcal{L}_{\mathrm{multi}},
  \label{eq:ch5:L1} \\
  \mathcal{L}^{(2)}
  &= \mathrm{CVaR}_{0.95}(J)
     + 0.3\,\mathcal{L}^{(1)}
     + w_\lambda(\|\lambda_s\|^2 + \|\lambda_h\|^2).
  \label{eq:ch5:L2}
\end{align}
Three structural changes are visible:
\begin{enumerate}[(i)]
  \item The primary loss switches from expected cost to tail-risk
        ($\mathrm{CVaR}_{0.95}$), because rare material hazards require
        a coherent risk measure, not just expectation minimization.
  \item The Stage~1 objective is \emph{retained} at weight $0.3$ as
        a scaffold, preventing the geometry backbone from degrading
        when the material heads are first unfrozen.
  \item Two new regularization terms $\|\lambda_s\|^2 + \|\lambda_h\|^2$
        prevent material forces from overpowering geometry prematurely.
\end{enumerate}

\begin{proposition}[CVaR subsumes expected cost]
\label{prop:ch5:cvar}
For any rollout cost distribution $J$,
$\mathrm{CVaR}_\beta(J) \geq \mathbb{E}[J]$ for all $\beta \in (0,1)$.
Moreover, $\mathrm{CVaR}_\beta(J) \to \mathrm{ess\,sup}(J)$ as
$\beta \to 1$.
A policy minimizing $\mathrm{CVaR}_{0.95}(J)$ therefore provides
stronger guarantees against catastrophic outcomes than a policy
minimizing $\mathbb{E}[J]$ alone.
\end{proposition}

\begin{proof}
Standard result from coherent risk measure theory~\citep{rockafellar2002cvar}.
The key property is that CVaR at level $\beta$ equals the expected
value over the worst-$(1-\beta)$ fraction of outcomes.
\end{proof}

\subsection{Enlargement D: Adaptation Timescales}
\label{subsec:ch5:timescalesD}

The most structurally significant enlargement is in the update hierarchy.
Stage~1 has a single online adaptation loop (segment-level barrier weight
updates).
Stage~2 has three nested loops:

\begin{equation}
  \underbrace{\Delta\theta_k \leftarrow
    (J_k^\top J_k + \lambda_J I)^{-1}J_k^\top\Delta\vec{o}_k}_{\text{segment (fast, local)}}
  \;\subset\;
  \underbrace{\theta \leftarrow \theta -
    \eta\nabla_\theta\mathcal{L}_{\mathrm{total}}}_{\text{episode (medium)}}
  \;\subset\;
  \underbrace{\Pi_{\mathrm{curr}}^{e+1} \propto
    1-\hat{c}_e}_{\text{curriculum (slow, global)}}.
  \label{eq:ch5:nested_loops}
\end{equation}

\begin{remark}[Why three loops are necessary]
\label{rem:ch5:loops}
Each loop corrects a different failure mode:
\begin{itemize}
  \item The \emph{segment loop} corrects local coefficient misalignment
        within a single corridor—fast enough to handle rapid context
        changes.
  \item The \emph{episode loop} optimizes the CVaR objective over the
        full trajectory—it needs the whole rollout to estimate tail risk.
  \item The \emph{curriculum loop} governs which configurations have
        been \emph{conquered}—it requires many episodes to assess
        stability and coverage.
\end{itemize}
Collapsing all three into a single gradient step would either be too
slow for local alignment or too noisy for global curriculum design.
The three loops are not an implementation convenience; they are
structurally necessary for the different timescales of the problem.
\end{remark}

% -----------------------------------------------------------------------
\section{What the Enrichments Buy: A Qualitative Analysis}
\label{sec:ch5:qualitative}
% -----------------------------------------------------------------------

\subsection{Failure Modes Addressed by Each Stage}
\label{subsec:ch5:failure}

\begin{table}[t]
\centering
\caption{%
  \textbf{Failure modes and which stage addresses them.}
  Each stage eliminates a specific class of failure and reveals the
  next-stage failure mode.
}
\label{tab:ch5:failures}
\small
\setlength{\tabcolsep}{5pt}
\renewcommand{\arraystretch}{1.3}
\begin{tabular}{lll}
\toprule
\textbf{Failure mode} & \textbf{Stage that exposes it} &
\textbf{Stage that addresses it} \\
\midrule
Oscillation / barrier ringing & Pre-Stage 1 (no dissipation) &
Stage 1 ($\mathcal{L}_{\mathrm{fric}}$ + explicit $R_\theta$) \\
Geometry-blind hazard entry & Stage 1 (no material channel) &
Stage 2 ($F_{\mathrm{soft}}, F_{\mathrm{hard}}$) \\
Rare catastrophic encounters & Stage 1 (expected-cost only) &
Stage 2 ($\mathrm{CVaR}_{0.95}$) \\
Local competence, global drift & Stage 1 (single timescale) &
Stage 2 (3-loop curriculum) \\
Oracle material dependence & Stage 2 (given $\tilde{r}, \phi$) &
Stage 3 (future: learned $\hat\rho$) \\
\bottomrule
\end{tabular}
\end{table}

Table~\ref{tab:ch5:failures} reveals a structural pattern: each stage
solves the failure modes exposed by the previous stage, but reveals
new ones in doing so.
This is the expected behavior of an enrichment hierarchy: not a
sequence of patches, but a series of genuinely harder subproblems.

\subsection{The Geometry Backbone as Invariant}
\label{subsec:ch5:invariant}

An important qualitative finding is that the geometry backbone
$H_{\mathrm{geom}}$ is preserved—not degraded—by the material-aware
enrichment.
On test cases where $\tilde{r} \approx 0$ and $\phi \gg 0$ (no material
risk), Stage~2 achieves path length and SPL within 3\% of Stage~1.
This confirms the backward-compatibility property
(Definition~\ref{def:ch5:compat}): the new coefficients default to near-zero
and leave the geometry backbone undisturbed.

The implication for the thesis program is:
\emph{each stage of enrichment is a stable foundation for the next},
not a replacement.
The enrichment ladder (Figure~\ref{fig:ch5:enrichment_ladder}) is not
a sequence of rewrites; it is an ascending structure where earlier stages
live inside later ones.

\subsection{Passivity as the Invariant of Enrichment}
\label{subsec:ch5:passivity}

The passivity condition $\mathcal{L}_{\mathrm{pass}} \leq \epsilon$
serves as the invariant that must be maintained across each enrichment.

\begin{proposition}[Passivity preservation under enrichment]
\label{prop:ch5:passivity_preservation}
Suppose $\mathfrak{S}^{(k)}$ satisfies the discrete passivity condition
$\Delta H^{(k)}_t \leq \tau y_t^\top u_t$ at convergence.
Then after enrichment to $\mathfrak{S}^{(k+1)}$ with the new energy
term $H_{\mathrm{new}}$ initialized at $\vartheta_{\mathrm{new}} = 0$
(backward-compatible initialization), the passivity condition is
satisfied at initialization:
$\Delta H^{(k+1)}_t\big|_{\vartheta_{\mathrm{new}}=0}
= \Delta H^{(k)}_t \leq \tau y_t^\top u_t$.
\end{proposition}

\begin{proof}
At $\vartheta_{\mathrm{new}} = 0$, $H^{(k+1)} = H^{(k)}$ exactly,
so $\Delta H^{(k+1)} = \Delta H^{(k)}$, which satisfies the passivity
condition by assumption.
\end{proof}

Proposition~\ref{prop:ch5:passivity_preservation} guarantees that
backward-compatible warm starts never \emph{introduce} passivity
violations at initialization.
Violations may appear as $\vartheta_{\mathrm{new}}$ is learned away
from zero; this is why $\mathcal{L}_{\mathrm{pass}}$ is monitored
throughout Phase~2 training and governs curriculum advancement.

% -----------------------------------------------------------------------
\section{Current Limitations of the Hierarchy}
\label{sec:ch5:limits}
% -----------------------------------------------------------------------

The enrichment principle and the two instantiated stages represent a
principled framework, but several limitations remain—some fundamental,
some practical.

\paragraph{Oracle material access.}
Stage~2 assumes that $\tilde{r}$ and $\phi$ are provided by an oracle.
In real deployment, these must be inferred from raw sensor data (RGB-IR,
multispectral, LiDAR reflection intensity).
Coupling perception co-training to the Hamiltonian learning loop is
Stage~3 of the enrichment ladder and is beyond the current scope.

\paragraph{Isotropic geometry metric.}
The mass matrix $\Mass$ in both stages is fixed (or mildly learned as
a scalar).
A richer geometry would use a learned anisotropic Riemannian metric
$\Mass(\Pos) \in \mathbb{S}^n_{++}$ that makes the kinetic energy
sensitive to local geometric curvature, gap width, and obstacle density.
This would improve navigation through highly anisotropic corridors but
requires a more complex gradient recursion.

\paragraph{Empirical coverage, not geometric coverage theory.}
The curriculum coverage criterion \eqref{eq:ch4:coverage} is empirical:
it counts conquered start–goal pairs from a finite sample.
A deeper theory would provide bounds on the sample complexity of
$\epsilon$-covering the configuration space, connecting curriculum
design to information-geometric arguments about Riemannian volume.

\paragraph{Stability guarantees are monitoring-based, not formal.}
The passivity criterion $\mathcal{L}_{\mathrm{pass}} \leq \epsilon$
is a sufficient heuristic for stability monitoring, not a formal
Lyapunov guarantee with a basin-of-attraction certificate.
Obtaining the latter would require bounding the Lipschitz constant of
the learned force field along the rollout, which depends on the
network architecture and is currently open.

\paragraph{Adaptation speed at the segment level.}
The segment Gauss-Newton update \eqref{eq:ch4:segment_update} depends
on a local Jacobian estimate that may be inaccurate in highly nonlinear
regions.
Trust-region or quasi-Newton variants could improve robustness, at the
cost of additional computation per segment.

% -----------------------------------------------------------------------
\section{Toward the Next Stage: Stage 3 and Beyond}
\label{sec:ch5:future}
% -----------------------------------------------------------------------

The enrichment principle predicts exactly what Stage~3 would require.
Replacing oracle material access with a learned belief map
$\hat\rho : \mathcal{Q}_f \to [0,1]$ (inferred from RGB-IR patches)
triggers the chain of consequences:
\begin{equation}
  H_{\mathrm{perc}}(\Pos;\hat\rho)
  \;\to\; F_{\mathrm{perc}} = -\nabla_\Pos H_{\mathrm{perc}}
  \;\to\; \vartheta_{\mathrm{perc}} = (\lambda_\rho)
  \;\to\; \mathcal{L}_{\mathrm{perc}}
  \;\to\; \text{belief update law}.
  \label{eq:ch5:stage3_chain}
\end{equation}
The belief update law would operate at a new timescale: fast enough to
track material transitions between segments, but slow enough to average
out sensor noise.
This is a fourth loop added to the update hierarchy $\mathcal{U}$.

Beyond Stage~3, the framework supports further enrichments:

\begin{itemize}
  \item \textbf{Anisotropic metric.}
        Adding $H_{\mathrm{metric}}(\Pos) = \frac{1}{2}\Mom^\top
        \Mass(\Pos)^{-1}\Mom$ with a position-dependent mass matrix
        makes kinetic energy sensitive to terrain curvature.
  \item \textbf{Multi-agent coordination.}
        A shared energy term $H_{\mathrm{social}}$ encoding inter-agent
        distances could extend the framework to cooperative navigation,
        where each agent's Hamiltonian couples to others'.
  \item \textbf{Topological coverage.}
        Using persistent homology of the visited configuration set as
        a coverage metric would replace empirical conquest counts with
        a topologically-grounded measure of how much of the space has
        been stably navigated.
\end{itemize}

Each of these would follow the same enrichment chain \eqref{eq:ch5:chain},
confirming that the framework is not a special-purpose system but a
general structured learner.

% % -----------------------------------------------------------------------
% \section{Chapter Summary}
% \label{sec:ch5:summary}
% % -----------------------------------------------------------------------

% \begin{figure}[htbp]
% \centering
% \begin{tikzpicture}[
%   font=\small\sffamily,
%   box/.style={rectangle, rounded corners=5pt,
%               draw=#1!70!black, fill=#1!10,
%               line width=1.1pt,
%               minimum width=3.6cm, minimum height=1.1cm,
%               align=center, inner sep=6pt},
%   arr/.style={-{Stealth[length=6pt]}, line width=1.2pt, color=#1},
%   lbl/.style={font=\scriptsize\sffamily\itshape, text=gray!60!black},
% ]
% % Row 1: skeleton
% \node[box=gray] (skel) at (0,0)
%   {\textbf{Shared skeleton}\\$\Phi_{\vartheta_t}(\PhaseState_t,y_t)$};

% % Row 2: two instances
% \node[box=teal, below left=1.2cm and 2.4cm of skel] (s1)
%   {\textbf{Stage 1}\\$H^{(1)},\,\mathcal{L}^{(1)},\,\mathcal{U}^{(1)}$};
% \node[box=orange, below right=1.2cm and 2.4cm of skel] (s2)
%   {\textbf{Stage 2}\\$H^{(2)},\,\mathcal{L}^{(2)},\,\mathcal{U}^{(2)}$};

% % Enrichment principle arrow
% \node[box=violet, below=2.8cm of skel] (ep)
%   {\textbf{Enrichment principle}\\Theorem~\ref{thm:ch5:enrichment}};

% % Row 3: thesis claim
% \node[box=gray!60!black,
%       minimum width=8.6cm,
%       below=1.4cm of ep] (claim)
%   {\textbf{Thesis claim:} structured progression, not two separate systems.\\
%    \textit{Each enrichment is a necessity, not a design choice.}};

% \draw[arr=gray!60]  (skel.south west) -- (s1.north);
% \draw[arr=gray!60]  (skel.south east) -- (s2.north);
% \draw[arr=teal]     (s1.south east)   -- (ep.north west);
% \draw[arr=orange]   (s2.south west)   -- (ep.north east);
% \draw[arr=violet, line width=1.4pt]  (ep.south)   -- (claim.north);

% \end{tikzpicture}
% \caption{%
%   \textbf{Chapter~\ref{chap:hierarchy} thesis diagram.}
%   A shared skeleton underlies both stages.
%   The enrichment principle (Theorem~\ref{thm:ch5:enrichment}) explains
%   why the transition from Stage~1 to Stage~2 is structurally forced,
%   not an ad-hoc addition.
%   The thesis claim is that this principled progression—not any single
%   stage—is the central contribution.
% }
% \label{fig:ch5:summary}
% \end{figure}

% This chapter synthesized Chapters~\ref{chap:grlsnam} and~\ref{chap:material}
% through four contributions:

% \begin{enumerate}
%   \item \textbf{The shared skeleton} (Definition~\ref{def:ch5:skeleton},
%         \eqref{eq:ch5:template}).
%         A single template $\Phi_{\vartheta_t}(\PhaseState_t,y_t)$
%         underlies both stages; what varies is the contents of
%         $H_\theta$, $R_\theta$, $\mathcal{L}$, and $\mathcal{U}$.

%   \item \textbf{The enrichment principle}
%         (Theorem~\ref{thm:ch5:enrichment}, \eqref{eq:ch5:chain}).
%         Adding one energy term forces five downstream consequences:
%         a new force channel, coefficient, objective term, gradient
%         pathway, and update law.
%         This chain makes the transition from Stage~1 to Stage~2 a
%         structural necessity, not an implementation choice.

%   \item \textbf{Backward compatibility and passivity as invariants}
%         (Definitions~\ref{def:ch5:compat} and
%         Proposition~\ref{prop:ch5:passivity_preservation}).
%         Each enrichment is safe: backward-compatible initialization
%         preserves passivity at initialization, and the curriculum
%         monitors it throughout.

%   \item \textbf{A clear path to Stage~3}
%         (\S\ref{sec:ch5:future}).
%         Replacing oracle material access with learned perception
%         follows the same enrichment chain, and the framework is
%         identified as a general platform for further enlargement.
% \end{enumerate}

% The central thesis claim, visualized in Figure~\ref{fig:ch5:summary},
% is that the structured progression from Stage~1 to Stage~2—and beyond—
% is the real contribution.
% Not a better navigator in any single setting, but a principled
% framework for building richer navigators by enriching a structured
% phase-space learner.